\scanpage{69}
\mspage{2}
then $\mathbf{\Gamma}$ acts not only on \deleted{$X$} the set $X$, but also on the groupoid $\widehat\Pi_1(X)$, because of the fact that the latter admits a purely scheme-theoretic definition, not involving the topology of $\mathbb C$. Of course, if $\lambda:\deleted{x} x\to y$ in $\Pi_1$, then for $u\in\Gamma$, $u(\lambda)$ is not in general an arrow of $\Pi_1$ (from $u(x)$ to $u(y)$) but only of $\widehat\Pi_1$.

\deleted{Notation} Notation ${}^ux$ ($x\in X,u\in\mathbf{\Gamma}$) and \deleted{\unsure{``}}
\[
{}^u\lambda:{}^ux-{}^uy\qquad\text{if }\lambda:x\to y\quad(\text{in }\widehat\Pi_1(X)).
\]
We have
\[
(6)\qquad\begin{cases}
{}^u(\operatorname{id}_x)=\operatorname{id}_{{}^ux}\qquad(x\in X)\\
{}^u(\mu\lambda)={}^u\mu\,{}^u\lambda
\end{cases}
\]
\diagram{\(\left(\begin{tikzcd}[ampersand replacement=\&,column sep=large]
x\arrow[r,"\lambda"]\&y\arrow[r,"\mu"]\&z\\
{}^ux\arrow[r,"{}^u\lambda"]\&{}^uy\arrow[r,"{}^u\mu"]\&{}^uz
\end{tikzcd}\right)\)}
Our object is to try to ``determine'' as far as possible this action of $\mathbf{\Gamma}$. To tell the truth, if $\bar K$ is the algebraic closure of $K$ in $\mathbb C$, we shall be interested above all in the subgroupoids of $\Pi_1(X)$ and $\widehat\Pi_1(X)$ induced on the set of objects $X_{\bar K}$, which we shall denote by $\Pi_1(X_{\bar K}),\widehat\Pi_1(X_{\bar K})$ -- the second is stable under the action of $\mathbf{\Gamma}$, and if we set
\[
(7)\qquad\Gamma=\operatorname{Aut}_{K\text{-cor}}(\bar K)=\operatorname{Gal}_{\bar K/K},
\]
so that $\Gamma$ is a quotient of $\mathbf{\Gamma}$, the action of $\mathbf{\Gamma}$ on $\Pi_1(X_{\bar K})$ factors through $\Gamma$. $\uparrow$
\margin{68}
